\begin{figure}[t]
\centering
\begin{tikzpicture}[node distance=14mm and 38mm]
  \node[screen] (home) {/};
  \node[screen, right=of home] (settings) {/settings};
  \node[screen, above=of home] (login) {/login};
  \node[screen, below=of home] (article) {/article/:slug};

  \draw[mustedge] (home) -- node[lbl, left] {\texttt{click} $\cdot$ must} (login);
  \draw[mayedge] (home.15) -- node[lbl, above=1mm] {\texttt{click} [\textit{requireAuth}] $\cdot$ may} (settings.165);
  \draw[mustedge] (settings.195) -- node[lbl, below=1mm] {\texttt{navigate} $\cdot$ must} (home.345);
  \draw[mayedge] (home) -- node[lbl, left] {\texttt{click} $\cdot$ may} (article);
\end{tikzpicture}
\caption{A small UI transition graph in the model of this thesis, with must and may edges.}
\label{fig:utg-example}
\par\smallskip{\footnotesize\noindent\emph{Note.} Each edge carries an event, an optional guard in brackets, and a modality: \emph{must} always fires when its code runs, \emph{may} is conditional (here on a route guard or a partly known target). In the measured graphs the unknown modality marks only edges into a sink for destinations that only the run time knows.}
\end{figure}
